\chapter{The Discrete Fourier Transform}

So far our Fourier story has lived in two slightly idealized worlds. In the
Fourier series we took a signal that repeats forever and broke it into a
\emph{discrete} list of harmonics. In the Fourier transform we took a signal of
finite energy that lives on the whole continuous line and broke it into a
\emph{continuous} smear of frequencies. Both are beautiful. Neither is what a
computer can actually hold.

A computer holds a finite list of numbers. Your microphone gives you, say,
$1024$ samples. Your IMU gives you a block of $256$ accelerometer readings. There
is no ``$t \to \infty$'', there is no integral over all of $\R$, there is just a
finite array. The Discrete Fourier Transform (the DFT) is the Fourier idea
rebuilt \emph{exactly} for that situation: finitely many samples in, finitely
many frequency numbers out, no approximation hidden anywhere, computable with a
finite number of multiplies and adds.

This is the version of Fourier you will actually run. Every spectrum analyzer,
every audio plugin, every gesture-recognition front end ultimately calls a DFT
(usually through its fast cousin, the FFT, which is the subject of the next
chapter). So we are going to build it slowly and geometrically, see it as a
matrix, understand exactly what each output number \emph{means}, and examine two
consequences that catch almost everyone off guard at first: the DFT implicitly treats
your $N$-sample block as one period of an infinitely repeating sequence, and any mismatch
at the block boundary causes \emph{spectral leakage}.

\begin{focusbox}
The DFT takes $N$ samples and reports how much
of each of $N$ evenly spaced rotating phasors is present --- and it
does this by taking inner products, which is orthogonal projection from linear
algebra carried out over $\C^N$.
\end{focusbox}

\section{From the continuous transform to a finite sum}
\label{ch08:from-continuous}

Recall the (continuous-time) Fourier transform of a signal $x(t)$:
\[
X(f) \;=\; \int_{-\infty}^{\infty} x(t)\, e^{-\jj 2\pi f t}\,\dd t .
\]
Geometrically this is one giant inner product: it slides a complex exponential
$e^{-\jj 2\pi f t}$ (a unit phasor spinning at frequency $f$) against $x(t)$ and
asks ``how much do you two overlap?''. We are going to discretize this idea in
two steps, and watch each idealization fall away.

\textbf{Step 1: we only have samples.} Suppose we sampled $x(t)$ every $T_s$
seconds, giving the sequence $x[n] = x(nT_s)$ for $n = 0, 1, \dots, N-1$. The
integral becomes a sum (this is just a Riemann sum, the bridge the calculus book
makes precise):
\[
X(f) \;\approx\; \sum_{n=0}^{N-1} x[n]\, e^{-\jj 2\pi f n T_s}\, T_s .
\]
We will drop the constant $T_s$ out front for now (it is just an overall scale
and we will track scaling carefully later). What remains depends on $f$ only
through the product $f T_s$, so define the \emph{normalized} angular frequency
$\omega = 2\pi f T_s$ (radians per sample). Then
\[
X(\omega) \;=\; \sum_{n=0}^{N-1} x[n]\, e^{-\jj \omega n}.
\]
This is the Discrete-Time Fourier Transform (DTFT). It eats a finite (or
infinite) sequence and produces a function of a \emph{continuous} variable
$\omega$. Better --- but still not finite output. A computer cannot store a
function of a continuous knob.

\textbf{Step 2: we only want finitely many frequencies.} Here is the clean move.
We have $N$ input numbers. Linear algebra tells us $N$ numbers carry exactly $N$
degrees of freedom, so we should ask for exactly $N$ frequency numbers back ---
no more (that would be redundant) and no fewer (that would lose information).
Which $N$ frequencies? We pick them \emph{evenly spaced} around the full circle of
possible $\omega \in [0, 2\pi)$:
\[
\omega_k \;=\; \frac{2\pi k}{N}, \qquad k = 0, 1, \dots, N-1.
\]
Sampling the DTFT at exactly these $N$ points gives the DFT:
\[
\boxed{\;X[k] \;=\; \sum_{n=0}^{N-1} x[n]\, e^{-\jj 2\pi k n / N}, \qquad k = 0,\dots,N-1.\;}
\]

\begin{intuitionbox}
Read the formula as a question repeated $N$ times. For each $k$ we cook up a
phasor that completes \emph{exactly $k$ full turns} over the $N$ samples, spin it
backwards against the data, add everything up, and the size of the result tells us
how strongly that particular ``$k$-turns-per-block'' rhythm is present. $k=0$ is
no turning at all (the average / DC level). $k=1$ is one slow turn across the
whole block. $k=N-1$ is the fastest turn the block can resolve.
\end{intuitionbox}

\begin{historybox}
The mathematics of interpolating periodic data from finite discrete points stretches back to
\textbf{Alexis Clairaut} (1754), \textbf{Joseph-Louis Lagrange} (1759), and \textbf{Carl Friedrich
Gauss} (1805), who developed trigonometric interpolation to calculate the orbits of asteroids.
However, the Discrete Fourier Transform as a standalone foundational pillar of digital engineering
only emerged with the rise of digital computers in the 1950s and 1960s.

When engineers first attempted to compute spectra numerically, they collided with two deep
conceptual traps. The first was \emph{periodic extension and spectral leakage}: the DFT is algebraically blind to the
world outside its $N$-sample window, implicitly assuming that the block $x[0],\dots,x[N-1]$
repeats periodically forever ($x[n+N] = x[n]$). If an incoming sinusoid does not complete an
exact integer number of cycles inside the window, the endpoints $x[0]$ and $x[N-1]$ do not match, creating an artificial discontinuity at every block seam. In the
frequency domain, this jump splatters energy across every frequency bin---an
artifact known as \emph{spectral leakage} that early computer analysts frequently misidentified
as physical overtones in their machinery.

The second trap was \emph{circular convolution}: in continuous calculus, multiplying two Fourier
transforms corresponds to linear convolution, but in the discrete domain, multiplying two $N$-point
DFTs $Y[k] = X[k]\cdot H[k]$ and taking the inverse DFT computes \emph{circular convolution} ($x[n] \circledast h[n]$),
where the trailing tail of the filtered response wraps around the buffer and corrupts the earliest
samples. Only by zero-padding the inputs to length at least $L+M-1$ can circular convolution be
forced to match linear filtering.
The DFT established digital frequency analysis, but evaluating it directly required an $N\times N$
matrix multiplication: $\mathcal{O}(N^2)$ operations. For large datasets, this computational wall
choked early mainframe computers until Cooley and Tukey's 1965 breakthrough, which forms the core of the
next chapter.
\end{historybox}

\section{The twiddle factor and the roots of unity}
\label{ch08:twiddle}

The repeated quantity $e^{-\jj 2\pi / N}$ shows up everywhere, so it gets a name.
Define the \emph{twiddle factor}
\[
W_N \;=\; e^{-\jj 2\pi / N}.
\]
Then the DFT and its inverse (which we will justify in
Section~\ref{ch08:orthogonality}) read compactly:
\[
X[k] = \sum_{n=0}^{N-1} x[n]\, W_N^{\,kn}, \qquad
x[n] = \frac{1}{N}\sum_{k=0}^{N-1} X[k]\, W_N^{\,-kn}.
\]
Geometrically $W_N$ is the point you land on after rotating clockwise by one
$N$-th of a full circle. Its powers $W_N^0, W_N^1, \dots, W_N^{N-1}$ are the $N$th
\emph{roots of unity}: $N$ points spread perfectly evenly around the unit circle.
That uniform spacing around the circle is what makes the DFT basis orthogonal.

\begin{figure}[H]
\centering
\begin{tikzpicture}[scale=2.1]
  \draw[->, >=Stealth, gray!50] (-1.3,0) -- (1.35,0) node[right, black]{\footnotesize Re};
  \draw[->, >=Stealth, gray!50] (0,-1.3) -- (0,1.35) node[above, black]{\footnotesize Im};
  \draw[thick,sourceblue] (0,0) circle (1);
  \foreach \k/\ang/\pos in {
    0/0/above right,
    1/-45/below right,
    2/-90/below right,
    3/-135/below left,
    4/-180/above left,
    5/-225/above left,
    6/-270/above right,
    7/-315/above right%
  } {
    \draw[gray!40,dashed] (0,0) -- (\ang:1);
    \filldraw[specred] (\ang:1) circle (0.035);
    \node[font=\footnotesize, \pos, outer sep=3pt] at (\ang:1) {$W_8^{\k}$};
  }
\end{tikzpicture}
\caption{The eight roots of unity $W_8^k = e^{-\jj 2\pi k/8}$ for $N=8$. The DFT
``measures'' a signal against phasors that sit at these evenly spaced angles. Note
the clockwise numbering, because of the minus sign in the exponent.}
\label{ch08:fig-roots}
\end{figure}

\begin{whybox}
Why \emph{evenly spaced} phasors and not some other set? Because evenly spaced
roots of unity are \emph{orthogonal}. Two of these sampled phasors, viewed as
vectors in $\C^N$, have zero inner product unless they are the same one. That is
precisely the property that made the Fourier series coefficients independent in
the earlier chapter, and it is what lets us invert the DFT cleanly. Orthogonality
is not a lucky accident; it is the reason this particular grid of frequencies is
the right grid. We prove it in Section~\ref{ch08:orthogonality}.
\end{whybox}

\section{The DFT as a matrix}
\label{ch08:matrix}

Here is the cleanest way to \emph{see} the DFT, and the place where the linear
algebra book pays off directly. The DFT is a linear map from $\C^N$ to $\C^N$, so
it \emph{is} a matrix. Stack the inputs into a vector $\vec{x} = (x[0], \dots,
x[N-1])^\top$ and the outputs into $\vec{X} = (X[0], \dots, X[N-1])^\top$. Then
\[
\vec{X} = \mat{F}_N\, \vec{x}, \qquad (\mat{F}_N)_{kn} = W_N^{\,kn} = e^{-\jj 2\pi k n / N}.
\]
Row $k$ of $\mat{F}_N$ is the conjugated phasor for frequency bin $k$, sampled at
the $N$ time points. Computing $X[k]$ is taking the dot product of that row with
your data. \emph{The DFT is matrix-vector multiplication}, full stop.

Let us write out $\mat{F}_4$ completely. With $N=4$, $W_4 = e^{-\jj\pi/2} = -\jj$,
so its powers cycle $W_4^0 = 1,\ W_4^1 = -\jj,\ W_4^2 = -1,\ W_4^3 = \jj$, then
repeat. The exponent is $kn \bmod 4$:
\[
\mat{F}_4 =
\begin{pmatrix}
W^0 & W^0 & W^0 & W^0\\
W^0 & W^1 & W^2 & W^3\\
W^0 & W^2 & W^4 & W^6\\
W^0 & W^3 & W^6 & W^9
\end{pmatrix}
=
\begin{pmatrix}
1 & 1 & 1 & 1\\
1 & -\jj & -1 & \jj\\
1 & -1 & 1 & -1\\
1 & \jj & -1 & -\jj
\end{pmatrix}.
\]

\begin{intuitionbox}
Look at the rows as ``test rhythms''. Row 0 is all ones: it just sums the data,
measuring the constant (DC) part. Row 2 alternates $1,-1,1,-1$: it measures the
fastest wiggle, the part that flips sign every sample (Nyquist). Rows 1 and 3
carry the genuine in-between rotations, and they are complex conjugates of each
other --- a fact that will become the symmetry of real-signal spectra.
\end{intuitionbox}

\begin{figure}[H]
\centering
\begin{tikzpicture}[x=0.92cm, y=1cm]
  % Column Headers
  \node[font=\small\bfseries, sigblue] at (2.25, 0.7) {Real Part: $\cos(2\pi k n / 8)$};
  \node[font=\small\bfseries, quotientorange] at (8.05, 0.7) {Imaginary Part: $\sin(2\pi k n / 8)$};

  \foreach \k/\ypos/\klabel in {
    0/0.0/{$k=0$ (DC, 0 cyc)},
    1/-1.35/{$k=1$ (1 cycle)},
    2/-2.70/{$k=2$ (2 cycles)},
    3/-4.05/{$k=3$ (3 cycles)},
    4/-5.40/{$k=4$ (Nyquist, 4 cyc)}
  } {
    % Row label on left
    \node[font=\footnotesize\bfseries, anchor=east] at (-0.25, \ypos) {\klabel};

    % Left baseline (Real)
    \draw[gray!35] (0, \ypos) -- (4.55, \ypos);
    % Right baseline (Imag)
    \draw[gray!35] (5.8, \ypos) -- (10.35, \ypos);

    % Faint underlying continuous sinusoids
    \draw[sigblue!35, smooth, domain=0:4.55, samples=100]
      plot (\x, {\ypos + 0.42*cos(deg(2*pi*\k*(\x/0.65)/8))});
    \draw[quotientorange!35, smooth, domain=0:4.55, samples=100]
      plot ({5.8 + \x}, {\ypos + 0.42*sin(deg(2*pi*\k*(\x/0.65)/8))});

    % Discrete samples at n=0..7
    \foreach \n in {0,...,7} {
      \pgfmathsetmacro{\xL}{\n*0.65}
      \pgfmathsetmacro{\xR}{5.8 + \n*0.65}
      \pgfmathsetmacro{\valRe}{0.42*cos(deg(2*pi*\k*\n/8))}
      \pgfmathsetmacro{\valIm}{0.42*sin(deg(2*pi*\k*\n/8))}

      % Real stem & dot
      \draw[sigblue, thick] (\xL, \ypos) -- (\xL, \ypos + \valRe);
      \fill[sigblue] (\xL, \ypos + \valRe) circle (1.6pt);

      % Imag stem & dot
      \draw[quotientorange, thick] (\xR, \ypos) -- (\xR, \ypos + \valIm);
      \fill[quotientorange] (\xR, \ypos + \valIm) circle (1.6pt);

      % Sample index tick on bottom row
      \ifnum\k=4
        \node[font=\tiny, text=gray!80!black, anchor=north] at (\xL, \ypos - 0.5) {$n{=}\n$};
        \node[font=\tiny, text=gray!80!black, anchor=north] at (\xR, \ypos - 0.5) {$n{=}\n$};
      \fi
    }
  }
\end{tikzpicture}
\caption{The DFT basis functions for $N=8$, showing the real (cosine, blue) and imaginary
(sine, orange) parts of the synthesis vectors $W_8^{-kn} = e^{\jj 2\pi kn/8}$ for $k=0,\dots,4$.
Each row is an orthogonal basis vector sampled at the 8 discrete time points $n=0,\dots,7$.
At $k=0$ (DC) the tone is constant and the imaginary part is zero; at $k=4$ (Nyquist) the
samples alternate $+1,-1,+1,-1$ at the maximum resolvable frequency; the higher bins
$k=5,6,7$ are conjugate mirrors of $k=3,2,1$ representing negative frequencies.}
\label{ch08:fig-basis}
\end{figure}

\begin{connectionbox}
This matrix view ties three books together. (1) Linear algebra: $\mat{F}_N$ is a
change of basis. The standard basis describes a signal by ``how much at each time
sample''; $\mat{F}_N$ rewrites it as ``how much at each frequency''. Same vector,
new coordinates. (2) The rows are orthogonal (next section), so $\mat{F}_N$ is ---
up to a scale --- a \emph{unitary} matrix, the complex cousin of a rotation. The
DFT does not stretch or distort the signal's geometry, it just rotates your
viewpoint. (3) Calculus: each entry samples the continuous kernel $e^{-\jj 2\pi f t}$ on a discrete
time--frequency grid.
\end{connectionbox}

\section{Why it inverts: orthogonality of the rows}
\label{ch08:orthogonality}

The single fact that makes the DFT useful is that you can get your signal
\emph{back}. Take the columns (equivalently rows, since $\mat{F}_N$ is symmetric)
and check their inner products. Consider two frequency vectors
$\vec{w}_k$ and $\vec{w}_m$ with entries $(\vec{w}_k)_n = W_N^{-kn}$ (the inverse
phasors). Their Hermitian inner product is
\[
\inner{\vec{w}_k}{\vec{w}_m}
= \sum_{n=0}^{N-1} W_N^{-kn}\, \overline{W_N^{-mn}}
= \sum_{n=0}^{N-1} e^{\jj 2\pi (k-m) n / N}.
\]
This is a finite geometric series in the ratio $r = e^{\jj 2\pi (k-m)/N}$. Two
cases:
\begin{itemize}
\item If $k = m$, every term is $e^0 = 1$, so the sum is $N$.
\item If $k \neq m$ (with $0 \le k,m < N$), then $r \neq 1$ but $r^N =
e^{\jj 2\pi(k-m)} = 1$, so $\sum_{n=0}^{N-1} r^n = \frac{r^N - 1}{r - 1} =
\frac{1-1}{r-1} = 0$.
\end{itemize}
Compactly,
\[
\sum_{n=0}^{N-1} e^{\jj 2\pi (k-m) n / N} = N\,\delta[k-m],
\]
where $\delta$ is $1$ at zero and $0$ otherwise. The frequency vectors are
mutually orthogonal, each with squared length $N$.

\begin{whybox}
That single identity \emph{is} the inverse DFT. Multiply the synthesis formula
$x[n] = \frac{1}{N}\sum_k X[k] W_N^{-kn}$ through: substitute the analysis formula
for $X[k]$, swap the order of summation, and the inner sum collapses to
$N\,\delta$ by the identity above, leaving exactly $x[n]$. The factor $1/N$ is
just the bookkeeping that undoes the squared-length $N$. Try this manipulation
yourself in Exercise~\ref{ch08:ex-prove-inverse}; it is three lines and it
demystifies the whole transform.
\end{whybox}

In matrix language the same fact reads $\mat{F}_N^{*} \mat{F}_N = N\,\mat{I}$,
where $\mat{F}_N^{*}$ is the conjugate transpose. Hence
\[
\mat{F}_N^{-1} = \tfrac{1}{N}\,\mat{F}_N^{*},
\]
which is why the inverse DFT is just the forward DFT with the sign of the exponent
flipped and a $1/N$ scale. Forward and inverse are nearly the same machine.

\begin{trapbox}
\textbf{Where does the $1/N$ go?} Conventions differ and this bites everyone.
Three common scalings: (a) $1/N$ on the inverse only (what we use, the default in
NumPy's \texttt{fft}/\texttt{ifft} and MATLAB); (b) $1/\sqrt{N}$ on \emph{both},
making $\mat{F}_N/\sqrt{N}$ exactly unitary and Parseval's theorem prettiest; (c)
$1/N$ on the forward only (common in some physics texts). They differ only by an
overall constant. Always check your library's convention before trusting absolute
magnitudes, especially when you compute power.
\end{trapbox}

\section{Reading the output: bins, resolution, and units}
\label{ch08:reading}

You ran an $N$-point DFT on data sampled at rate $f_s = 1/T_s$. You got back $N$
complex numbers $X[0], \dots, X[N-1]$. What do they \emph{mean} physically?

\textbf{Each output is a bin.} Output $X[k]$ corresponds to the analog frequency
\[
f_k = \frac{k}{N}\, f_s = \frac{k f_s}{N} \quad\text{Hz} \qquad (k = 0,\dots,N-1).
\]
The bins are equally spaced by
\[
\Delta f = \frac{f_s}{N} = \frac{1}{N T_s} = \frac{1}{\text{record length in seconds}}.
\]
This $\Delta f$ is the \emph{frequency resolution}: the smallest frequency
difference the transform can place in separate bins.

\begin{happeningbox}
The resolution formula $\Delta f = 1/(NT_s)$ carries a sobering consequence:
\textbf{resolution is set by how long you watched, not by how fast you
sampled}. $NT_s$ is the total duration of your record. Watch for $1$ second and
your bins are $1$\,Hz apart, no matter whether you sampled at $1$\,kHz or
$1$\,MHz. To tell two close tones apart you must observe longer. Sampling faster
only buys you a higher \emph{top} frequency, not finer spacing. This is a baby
version of the uncertainty principle: you cannot be simultaneously sharp in time
(short record) and sharp in frequency (fine bins).
\end{happeningbox}

\textbf{Magnitude and phase.} Each $X[k]$ is complex; write $X[k] = \abs{X[k]}\,
e^{\jj\angle X[k]}$. The magnitude $\abs{X[k]}$ says \emph{how much} of that
frequency is present; the phase $\angle X[k]$ says \emph{where} that sinusoid sits
in time (its alignment). Most spectrum plots show magnitude (or power
$\abs{X[k]}^2$) because the eye reads ``which frequencies are loud'' most easily,
but phase carries the timing and is essential for reconstruction --- throw it away
and you cannot get the signal back.

\textbf{The upper half is a mirror (for real signals).} If $x[n]$ is real --- as
all physical sensor data is --- then
\[
X[N-k] = \overline{X[k]}.
\]
So the second half of the spectrum is the complex conjugate of the first half
read backwards. It carries no new information. The genuinely distinct frequencies
live in bins $0$ through $N/2$. Bins above $N/2$ are the ``negative frequencies''
folded around; bin $k$ for $k > N/2$ really represents $-(N-k)$ cycles, i.e. the
same physical frequency rotating the other way.

\begin{figure}[H]
\centering
\begin{tikzpicture}
\begin{axis}[
  width=12cm, height=5.2cm,
  xlabel={bin $k$}, ylabel={$|X[k]|$},
  xmin=-0.5, xmax=15.5, ymin=0, ymax=9,
  xtick={0,2,4,6,8,10,12,14},
  ymajorgrids, grid style=gray!25,
  axis lines=left,
  title style={font=\small}, title={Magnitude spectrum of a real $N=16$ signal},
  ybar, bar width=7pt,
]
% N=16, a tone at k=3 (and its mirror at 13), plus DC
\addplot[draw=sourceblue, fill=fillsig] coordinates {
 (0,1.2) (1,0.3) (2,0.5) (3,8) (4,0.6) (5,0.2) (6,0.15) (7,0.1)
 (8,0.1) (9,0.1) (10,0.15) (11,0.2) (12,0.6) (13,8) (14,0.5) (15,0.3)
};
\draw[densely dashed, specred] (axis cs:8,0) -- (axis cs:8,9);
\node[specred, font=\scriptsize, anchor=south] at (axis cs:8,8.2) {Nyquist};
\node[font=\scriptsize, anchor=south] at (axis cs:3,8) {$k=3$};
\node[font=\scriptsize, anchor=south] at (axis cs:13,8) {mirror};
\end{axis}
\end{tikzpicture}
\caption{Magnitude spectrum of a real $N=16$ signal containing a single tone at
$3$ cycles per block (plus a small DC offset). The peak at $k=3$ is mirrored at
$k=13=N-3$. Everything to the right of the Nyquist bin $k=8$ is redundant for real
data. With $f_s = 16$\,Hz this peak sits at $f_3 = 3\cdot 16/16 = 3$\,Hz.}
\label{ch08:fig-spectrum}
\end{figure}

\begin{example}[Putting numbers on the axes]
\label{ch08:ex-axes}
You record an IMU axis at $f_s = 200$\,Hz and collect $N = 256$ samples.
\begin{itemize}
\item Record length: $N T_s = 256/200 = 1.28$\,s.
\item Bin spacing: $\Delta f = 200/256 \approx 0.781$\,Hz.
\item Bin $k$ sits at $f_k = 0.781\,k$\,Hz. Bin $5$ is at $3.9$\,Hz.
\item Nyquist (bin $128$) is at $100$\,Hz; bins $129$--$255$ mirror bins
$127$--$1$.
\item A $4$\,Hz hand tremor lands near bin $4\,/\,0.781 \approx 5.1$ --- between
bins $5$ and $6$, which is exactly the situation that causes leakage
(Section~\ref{ch08:leakage}).
\end{itemize}
\end{example}

\section{Two worked DFTs by hand}
\label{ch08:byhand}

Nothing builds intuition like grinding a small transform by hand. We do two.

\begin{example}[A 4-point DFT of a real signal]
\label{ch08:ex-4pt}
Let $x = (1, 2, 3, 4)$. Using $\mat{F}_4$ from Section~\ref{ch08:matrix},
\[
\vec{X} = \mat{F}_4 \vec{x} =
\begin{pmatrix}
1 & 1 & 1 & 1\\
1 & -\jj & -1 & \jj\\
1 & -1 & 1 & -1\\
1 & \jj & -1 & -\jj
\end{pmatrix}
\begin{pmatrix} 1\\2\\3\\4 \end{pmatrix}.
\]
Row by row:
\begin{align*}
X[0] &= 1+2+3+4 = 10, \\
X[1] &= 1 - 2\jj - 3 + 4\jj = -2 + 2\jj, \\
X[2] &= 1 - 2 + 3 - 4 = -2, \\
X[3] &= 1 + 2\jj - 3 - 4\jj = -2 - 2\jj.
\end{align*}
So $\vec{X} = (10,\ -2+2\jj,\ -2,\ -2-2\jj)$. Sanity checks: $X[0]=10$ is the sum
(four times the average $2.5$). And $X[3] = \overline{X[1]}$, exactly the
real-signal mirror symmetry $X[N-k]=\overline{X[k]}$. The magnitudes are
$(10,\ 2\sqrt2,\ 2,\ 2\sqrt2)$.
\end{example}

\begin{example}[Recovering a pure tone exactly]
\label{ch08:ex-tone}
Now take $N=8$ and the sampled cosine $x[n] = \cos(2\pi \cdot 2 \cdot n / 8) =
\cos(\pi n/2)$, which completes exactly $2$ cycles in the block. The samples are
\[
x = (1,\ 0,\ -1,\ 0,\ 1,\ 0,\ -1,\ 0).
\]
Because the tone fits a \emph{whole} number of cycles in the window, the DFT is as
clean as it ever gets. Using Euler, $\cos\theta = \tfrac12 e^{\jj\theta} +
\tfrac12 e^{-\jj\theta}$, so the signal is half a phasor at $+2$ cycles and half a
phasor at $-2$ cycles. The DFT therefore puts energy in exactly two bins:
\[
X[2] = \tfrac{N}{2} = 4, \qquad X[6] = \tfrac{N}{2} = 4,
\]
and $X[k] = 0$ for all other $k$. (Bin $6 = N-2$ is the mirror, the $-2$-cycle
component.) Verify $X[2]$ directly:
\[
X[2] = \sum_{n=0}^{7} x[n]\, e^{-\jj 2\pi \cdot 2 n / 8}
= \sum_n x[n]\, e^{-\jj \pi n / 2} = \sum_n x[n]\,(\cos\tfrac{\pi n}{2} - \jj\sin\tfrac{\pi n}{2}).
\]
The $\sin$ part hits zero wherever $x[n]\neq 0$, and the $\cos$ part reproduces
$x[n]$ itself, so $X[2] = \sum_n x[n]^2 = 1+1+1+1 = 4$. A tone that fits the
window perfectly produces a single sharp spike. Remember that phrase ``fits the
window perfectly'' --- the moment it does not, leakage appears.
\end{example}

\begin{figure}[H]
\centering
\begin{tikzpicture}
\begin{axis}[
  width=13cm, height=4.2cm,
  xlabel={$n$}, ylabel={$x[n]$},
  xmin=-0.4, xmax=7.4, ymin=-1.4, ymax=1.4,
  xtick={0,1,2,3,4,5,6,7}, ytick={-1,0,1},
  axis lines=left, ymajorgrids, grid style=gray!20,
  title style={font=\small}, title={$N=8$ sampled sinusoid: exactly $2$ cycles in the window}
]
\addplot[gray!55, thick, domain=0:7, samples=120] {cos(deg(pi*x/2))};
\addplot[only marks, mark=*, mark size=2pt, sourceblue]
  coordinates {(0,1)(1,0)(2,-1)(3,0)(4,1)(5,0)(6,-1)(7,0)};
\end{axis}
\end{tikzpicture}
\caption{The eight samples of Example~\ref{ch08:ex-tone}: $x[n]=\cos(\pi n/2)$,
exactly $2$ cycles across the $N=8$ block (samples as dots, underlying cosine in
gray). Because a whole number of cycles fits the window, the periodic extension
closes without a seam and the whole tone lands in bins $2$ and $6$.}
\label{ch08:fig-tone-time}
\end{figure}

\begin{labbox}
Reproduce both examples in three lines of NumPy and confirm by eye:
\begin{lstlisting}[language=Python]
import numpy as np
print(np.fft.fft([1,2,3,4]))            # -> [10, -2+2j, -2, -2-2j]
n = np.arange(8)
print(np.round(np.fft.fft(np.cos(np.pi*n/2)).real, 3))  # spike at 2 and 6
\end{lstlisting}
Then break it on purpose: change the tone to $2.5$ cycles per block
(\texttt{np.cos(2*np.pi*2.5*n/8)}) and watch the single spike smear across many
bins. That smear is the subject of the next section.
\end{labbox}

\section{The hidden assumption: the DFT is circular}
\label{ch08:circular}

Here is the most important conceptual point in the chapter, and the one most
likely to surprise you. The DFT does not see your $N$ samples as a finite snippet
sitting in silence. It sees them as \emph{one period of an infinitely repeating
signal}. The block wraps around: sample $N$ would be sample $0$ again.

Why? Look at the basis phasors $e^{-\jj 2\pi k n/N}$. Every one of them satisfies
$e^{-\jj 2\pi k (n+N)/N} = e^{-\jj 2\pi k n/N}$ --- they are all periodic with
period $N$. If your toolbox is built entirely from period-$N$ pieces, then
\emph{anything} you build is period-$N$. The DFT literally cannot represent a
signal that is not periodic in $N$; it represents the periodic extension of your
data whether you wanted that or not.

\begin{figure}[H]
\centering
\begin{tikzpicture}
\begin{axis}[
  width=13cm, height=4.4cm,
  xlabel={$n$}, ylabel={$x[n]$},
  xmin=-0.5, xmax=23.5, ymin=-1.4, ymax=1.4,
  xtick={0,4,8,12,16,20}, ytick={-1,0,1},
  axis lines=left, ymajorgrids, grid style=gray!20,
  title style={font=\small}, title={Periodic extension of an off-bin sinusoid ($2.5$ cycles)}
]
% one period (n=0..7) of a tone that does NOT fit: 2.5 cycles, then repeat showing the jump
\addplot[only marks, mark=*, mark size=1.4pt, sourceblue, domain=0:7, samples=8]
   {cos(deg(2*pi*2.5*x/8))};
\addplot[only marks, mark=o, mark size=1.6pt, specred, domain=8:15, samples=8]
   {cos(deg(2*pi*2.5*(x-8)/8))};
\addplot[only marks, mark=o, mark size=1.6pt, specred, domain=16:23, samples=8]
   {cos(deg(2*pi*2.5*(x-16)/8))};
\draw[gray!55] (axis cs:7.5,-1.4) -- (axis cs:7.5,1.4);
\draw[gray!55] (axis cs:15.5,-1.4) -- (axis cs:15.5,1.4);
\draw[->,thick,specred] (axis cs:7,-0.71) to[bend left=40] (axis cs:8,0.71);
\node[specred,font=\scriptsize,anchor=south west] at (axis cs:7.4,0.0) {jump!};
\end{axis}
\end{tikzpicture}
\caption{The DFT's view of a $2.5$-cycle tone in an $8$-sample window: it tiles the
block end to end (filled dots = your data, open dots = the implied repeats). Since
$2.5$ cycles do not close smoothly, a discontinuity appears at every seam (red
arrow). The DFT must spend energy across many bins to manufacture that corner ---
that is leakage.}
\label{ch08:fig-circular}
\end{figure}

\begin{happeningbox}
This circularity is not a bug, it is the exact mathematical content of the DFT.
Everything ``circular'' in DSP flows from it: circular convolution (the DFT
multiplies spectra, which in time means wrap-around convolution, not the linear
convolution you usually want), circular shifts, the mirror symmetry of real
spectra, and leakage. Hold onto the mental image of your data \emph{taped into a
ring}. Many later confusions evaporate the moment you remember the DFT lives on a
circle, not a line.
\end{happeningbox}

\section{Spectral leakage}
\label{ch08:leakage}

Now we can explain the smear you saw in the lab box. When a tone completes a
\emph{whole} number of cycles in the window (Example~\ref{ch08:ex-tone}), the
periodic extension has no jump: the ring closes perfectly, the tone is genuinely
one of the basis phasors, and all its energy lands in one bin. This is called an
\emph{on-bin} or \emph{coherent} frequency.

But a real $4$\,Hz tremor sampled into a $1.28$\,s window (Example~\ref{ch08:ex-axes})
sits at bin $5.1$ --- \emph{between} bins. Its periodic extension has a jump at
every seam (Figure~\ref{ch08:fig-circular}). A jump is sharp, and sharp features
need many frequencies to build, so the DFT spreads the tone's energy across
\emph{all} the bins, strongest near $5$ and $6$ but with a slowly decaying skirt
everywhere. That spreading is \textbf{spectral leakage}: energy ``leaks'' out of
the true frequency into neighboring bins.

\begin{figure}[H]
\centering
\begin{tikzpicture}
\begin{axis}[
  width=12.5cm, height=5.2cm,
  xlabel={bin $k$}, ylabel={$|X[k]|$ (linear)},
  xmin=-0.5, xmax=16.5, ymin=0, ymax=9,
  xtick={0,2,4,6,8,10,12,14,16},
  axis lines=left, ymajorgrids, grid style=gray!25,
  title style={font=\small}, title={Spectral leakage: on-bin ($4$ cycles) vs off-bin ($4.5$ cycles)},
  ybar, bar width=5pt,
  area legend,
  legend style={font=\scriptsize, at={(0.98,0.97)}, anchor=north east},
]
% on-bin: clean spike at k=4
\addplot[draw=sourceblue, fill=fillsig] coordinates {
 (0,0.05)(1,0.05)(2,0.05)(3,0.05)(4,8)(5,0.05)(6,0.05)(7,0.05)
 (8,0.05)(9,0.05)(10,0.05)(11,0.05)(12,0.05)(13,0.05)(14,0.05)(15,0.05)(16,0.05)};
\addlegendentry{on-bin (4.0 cyc)}
% off-bin: smeared around 4.5
\addplot[draw=specred, fill=specred!35] coordinates {
 (0,0.4)(1,0.5)(2,0.7)(3,1.4)(4,4.6)(5,4.6)(6,1.4)(7,0.7)(8,0.5)
 (9,0.4)(10,0.35)(11,0.3)(12,0.3)(13,0.28)(14,0.27)(15,0.27)(16,0.27)};
\addlegendentry{off-bin (4.5 cyc)}
\end{axis}
\end{tikzpicture}
\caption{On-bin versus off-bin tones, $N=32$ (first half shown). The on-bin tone
(blue, exactly $4$ cycles per block) lands in a single bin. The off-bin tone (red,
$4.5$ cycles) splits between bins $4$ and $5$ and leaks a decaying skirt across
every other bin. The peak is also lower --- energy that scattered is energy not in
the main bin (this height loss is called \emph{scalloping}).}
\label{ch08:fig-leakage}
\end{figure}

\begin{whybox}
There is a sharper way to see leakage that connects straight to the continuous
transform of Chapter~7. Taking only $N$ samples is the same as multiplying an
infinitely long signal by a \emph{rectangular window} (1 inside the block, 0
outside). Multiplication in time is convolution in frequency. The Fourier
transform of a rectangle is a sinc, so the true line spectrum of a pure tone gets
\emph{convolved with a sinc}. When the tone is on-bin, the DFT samples that sinc
exactly at its zero crossings except for the one peak, so you see a clean spike.
When the tone is off-bin, the DFT samples the sinc \emph{between} its zeros and you
see the whole ringing shape. Leakage is the sinc's side-lobes, sampled. We will
make this precise when we study windowing and the convolution theorem.
\end{whybox}

\subsection*{Taming leakage: windowing (a preview)}

The rectangular window's side-lobes fall off slowly (only $\sim 13$\,dB down from
the peak). The fix is to multiply your $N$ samples by a smooth \emph{window}
function $w[n]$ that tapers gently to zero at both ends \emph{before} taking the
DFT. A tapered block has no jump at the seam, so the periodic extension is smooth,
so far less energy leaks. The price: the smooth taper widens the main lobe, so you
trade resolution for lower leakage. You cannot escape the trade-off, only choose
where to sit on it.

\begin{figure}[H]
\centering
\begin{tikzpicture}
\begin{axis}[
  width=11cm, height=4.4cm,
  xlabel={$n$}, ylabel={$w[n]$},
  xmin=0, xmax=31, ymin=-0.05, ymax=1.15,
  xtick={0,8,16,24,31}, ytick={0,0.5,1},
  axis lines=left, ymajorgrids, grid style=gray!20,
  title style={font=\small}, title={Comparison of window functions ($N=32$)},
  legend style={font=\scriptsize, at={(0.5,-0.30)}, anchor=north, legend columns=3},
]
\addplot[sourceblue, thick, domain=0:31, samples=32, mark=none] {1};
\addlegendentry{rectangular}
\addplot[specred, thick, domain=0:31, samples=64, mark=none]
  {0.5 - 0.5*cos(deg(2*pi*x/31))};
\addlegendentry{Hann}
\addplot[targetgreen, thick, dashed, domain=0:31, samples=64, mark=none]
  {0.54 - 0.46*cos(deg(2*pi*x/31))};
\addlegendentry{Hamming}
\end{axis}
\end{tikzpicture}
\caption{Three windows over $N=32$ samples. The rectangular window (do nothing)
keeps full amplitude but jumps at the edges; the Hann and Hamming windows taper to
(near) zero, smoothing the seam and suppressing leakage at the cost of a wider main
lobe. Chapter~21 develops the full menu and their trade-offs.}
\label{ch08:fig-windows}
\end{figure}

\begin{trapbox}
Do not confuse \emph{zero-padding} with better resolution. Appending zeros to your
block before the DFT (say from $N=256$ to $N=1024$) gives you more bins, so the
spectrum looks smoother and you can read peak locations more finely. But it adds
\emph{no new information}: the underlying record length $NT_s$ that sets $\Delta f$
did not change, you just interpolated the same DTFT more densely. Zero-padding
\emph{interpolates}; only observing longer truly \emph{resolves}.
\end{trapbox}

\section{Parseval: energy is preserved}
\label{ch08:parseval}

Because $\mat{F}_N/\sqrt{N}$ is unitary, the DFT preserves length (energy) up to
the scale factor. Concretely,
\[
\sum_{n=0}^{N-1} \abs{x[n]}^2 = \frac{1}{N}\sum_{k=0}^{N-1} \abs{X[k]}^2 .
\]
The total energy in the time samples equals the total energy across the bins
(divided by $N$). This is the discrete Parseval / Rayleigh identity, and it is the
same statement as ``a rotation does not change a vector's length'', now in $\C^N$.

\begin{example}[Check Parseval on the 4-point example]
\label{ch08:ex-parseval}
From Example~\ref{ch08:ex-4pt}, the time energy is $1^2+2^2+3^2+4^2 = 30$. The
spectral energy is $\abs{10}^2 + \abs{-2+2\jj}^2 + \abs{-2}^2 + \abs{-2-2\jj}^2 =
100 + 8 + 4 + 8 = 120$, and $120/N = 120/4 = 30$. They match.
\end{example}

\begin{connectionbox}
Parseval is why we can compute signal \emph{power} in the frequency domain ---
an indispensable foundation for spectral estimation (Chapter~21) and for Chapter~20. If a sensor
adds white noise of variance $\sigma^2$ per sample, that energy spreads evenly
across all $N$ bins, so each bin carries an expected $\sigma^2$ of noise power
while a coherent tone piles all its energy into one bin. That contrast --- noise
flat, signal concentrated --- is the entire reason the DFT is such an effective
detector. We will quantify the resulting ``processing gain'' in Chapter~20.
\end{connectionbox}

\section{Cost, and a look ahead to the FFT}
\label{ch08:cost}

The matrix $\mat{F}_N$ is $N\times N$ and dense, so computing $\vec{X} =
\mat{F}_N\vec{x}$ directly costs about $N^2$ complex multiply-adds. For $N=1024$
that is roughly a million operations \emph{per transform}; for the real-time
audio rate of dozens of transforms per second it quickly becomes a problem on a
small microcontroller.

The next chapter shows that $\mat{F}_N$ has enormous hidden structure --- the
twiddle factors repeat in a recursive pattern --- which lets us factor the matrix
into a product of very sparse matrices and slash the cost to $\sim N\log_2 N$.
That is the Fast Fourier Transform. Everything you understood here about bins,
leakage, and circularity carries over unchanged; the FFT computes \emph{exactly}
the same numbers, only faster.

\begin{goalbox}
For the wrist gesture controller, the DFT is the front-end that turns raw motion
into features. A practical pipeline: sample the IMU at $f_s = 200$\,Hz, buffer
$N=256$ samples ($1.28$\,s), apply a Hann window, take the DFT of each axis, and
read off band energies --- e.g. how much power sits in $1$--$5$\,Hz (deliberate
gesture motion) versus $8$--$12$\,Hz (physiological tremor) versus the near-DC bin
(gravity / orientation). Those band energies are compact, robust features for a
gesture classifier (Figure~\ref{ch08:fig-bands}).

\textbf{The limits, stated plainly.} A wrist-worn IMU alone senses the \emph{gross} motion of
the whole forearm and hand. Its spectrum can reliably separate coarse, discrete
gestures (a flick, a twist, a tap, ``holding still'') and tell rhythmic motion
from static pose. It \emph{cannot} recover individual finger joint angles --- the
fingers barely move the wrist, so their signal is essentially absent from a wrist
IMU. Continuous finger pose needs sensors on the fingers (instrumented rings or
gloves); finger-level \emph{activation} needs forearm EMG or force myography (FMG),
which read muscle electrical or mechanical activity. The DFT will be just as
central there --- EMG band power is a standard feature --- but be clear-eyed about
what each sensor can and cannot deliver. Promising ``per-finger MIDI'' from a lone
wrist IMU is a promise the physics will not keep.

\begin{figure}[H]
\centering
\begin{tikzpicture}
\begin{axis}[
  width=12.5cm, height=5.2cm,
  xlabel={frequency (Hz)}, ylabel={$|X[k]|$},
  xmin=0, xmax=20, ymin=0, ymax=6.8,
  xtick={0,2,4,6,8,10,12,14,16,18,20},
  ytick={0,2,4,6},
  axis lines=left, ymajorgrids, grid style=gray!25,
  title style={font=\small}, title={Windowed IMU magnitude spectrum and gesture feature bands},
  ybar, bar width=4pt,
]
\addplot[draw=sourceblue, fill=fillsig] coordinates {
 (0.78,5.2)(1.56,3.1)(2.34,4.2)(3.12,2.6)(3.91,1.7)(4.69,1.1)
 (5.47,0.5)(6.25,0.4)(7.03,0.35)(7.81,0.6)(8.59,1.9)(9.38,2.4)
 (10.16,2.0)(10.94,1.2)(11.72,0.5)(12.5,0.3)(13.28,0.25)(14.06,0.2)
 (14.84,0.2)(15.62,0.18)(16.41,0.15)(17.19,0.15)(18.0,0.12)(18.75,0.12)};
\draw[densely dashed, structurepurple] (axis cs:1,0) -- (axis cs:1,5.85);
\draw[densely dashed, structurepurple] (axis cs:5,0) -- (axis cs:5,5.85);
\draw[densely dashed, quotientorange] (axis cs:8,0) -- (axis cs:8,5.85);
\draw[densely dashed, quotientorange] (axis cs:12,0) -- (axis cs:12,5.85);
\node[structurepurple, font=\tiny, anchor=south] at (axis cs:3,4.85) {gesture 1--5 Hz};
\node[quotientorange, font=\tiny, anchor=south] at (axis cs:10,4.85) {tremor 8--12 Hz};
\node[specred, font=\tiny, anchor=south west] at (axis cs:0.35,6.05) {near-DC: gravity};
\draw[->, specred, thin] (axis cs:0.78,6.05) -- (axis cs:0.78,5.35);
\end{axis}
\end{tikzpicture}
\caption{Sketch of a windowed wrist-IMU magnitude spectrum during a slow gesture
($f_s=200$\,Hz, $N=256$, so $\Delta f \approx 0.78$\,Hz). Summing $|X[k]|^2$ over
the $1$--$5$\,Hz band captures deliberate motion; the $8$--$12$\,Hz band tracks
physiological tremor; the near-DC bin carries the gravity vector. These three band
energies make compact gesture features.}
\label{ch08:fig-bands}
\end{figure}
\end{goalbox}

\begin{buildbox}
A minimal fixed-block DFT loop in C, suitable as a reference before you drop in a
real FFT library (CMSIS-DSP on Cortex-M, or KissFFT on a Pi). Keep $N$ a power of
two so swapping in an FFT later is painless.
\begin{lstlisting}[language=C]
#include <math.h>
#define N 256
// in[] = windowed real samples; out_mag[] = magnitude of bins 0..N/2
void dft_mag(const float in[N], float out_mag[N/2 + 1]) {
    for (int k = 0; k <= N/2; ++k) {     // real signal: half the bins suffice
        float re = 0.f, im = 0.f;
        for (int n = 0; n < N; ++n) {
            float ang = -2.f * (float)M_PI * k * n / N;
            re += in[n] * cosf(ang);
            im += in[n] * sinf(ang);
        }
        out_mag[k] = sqrtf(re*re + im*im);
    }
}
\end{lstlisting}
This is $O(N^2)$ and recomputes cosines in the inner loop --- fine for
understanding and for tiny $N$, far too slow for $N=256$ at $200$\,Hz on an 8-bit
AVR. Precompute a twiddle table, or just use an FFT. The point of writing it once
by hand is so the FFT is never a black box.
\end{buildbox}

\begin{aestheticbox}
Look at how the continuous theory compressed into a finite computation. We started with an integral over all of
time against a continuum of frequencies --- an object no computer can touch --- and
by sampling \emph{both} time and frequency we landed on a single finite matrix
that is, up to scale, a rotation of $\C^N$. The same orthogonality that powered
Fourier series, the same projection idea from linear algebra, the same
time/frequency uncertainty from calculus, all reappear in finite, computable form.
The DFT is where the grand Fourier theory becomes something you can run on a chip
strapped to your wrist. That is a surprising amount of unity in one $N\times N$
array.
\end{aestheticbox}

\begin{connectionbox}
\textbf{Capabilities, Downsides, and Forward Roadmap}
\begin{itemize}
  \item \textbf{What this unlocked}: The DFT brings frequency analysis into the digital age: an exact,
  computable, finite-dimensional unitary matrix transformation $\vec{X} = \mat{F}_N \vec{x}$ that
  runs directly on digital hardware. Energy is preserved by Parseval's theorem, and signals are
  decomposed onto an orthonormal basis of discrete complex exponentials.
  \item \textbf{Honest downsides and limitations}:
  \begin{enumerate}
    \item \emph{Computational Complexity}: Direct matrix-vector evaluation costs $\mathcal{O}(N^2)$
    complex operations. For real-time embedded audio ($N = 2048$) or wearable processing, this
    quadratically explodes the processor cycle budget.
    \item \emph{Leakage vs.\ Resolution}: Frequencies falling between bin centers leak energy across
    the spectrum; applying tapered windows mitigates leakage but widens the main lobe, smearing
    closely spaced spectral peaks.
    \item \emph{Circular Wrap-Around}: Straightforward multiplication in the DFT domain executes
    circular rather than linear convolution unless padded with extra zeros.
  \end{enumerate}
  \item \textbf{Where we go next}: Chapter 9 solves the computational bottleneck with the Cooley--Tukey
  Fast Fourier Transform (FFT), reducing the runtime from $\mathcal{O}(N^2)$ to $\mathcal{O}(N\log_2 N)$.
  Chapter 10 overcomes the stationary block assumption with the Short-Time Fourier Transform (STFT);
  and Chapter 13 uses zero-padded FFTs to compute linear convolutions in $\mathcal{O}(N\log N)$ time.
\end{itemize}
\end{connectionbox}

\section{Exercises}
\label{ch08:exercises}

\begin{exercise}
\label{ch08:ex-prove-inverse}
Prove the inverse DFT. Start from $y[n] = \frac{1}{N}\sum_{k=0}^{N-1} X[k]
e^{\jj 2\pi kn/N}$, substitute $X[k] = \sum_{m} x[m] e^{-\jj 2\pi km/N}$, swap the
two sums, and apply the orthogonality identity
$\sum_{k} e^{\jj 2\pi k(n-m)/N} = N\,\delta[n-m]$ to show $y[n] = x[n]$. Identify
where the factor $1/N$ is consumed.
\end{exercise}

\begin{exercise}
Compute the $4$-point DFT of $x = (1,0,0,0)$ (a unit impulse) and of $x =
(1,1,1,1)$ (a constant). Explain in one sentence each why one spectrum is flat and
the other is a single spike. What general rule about ``narrow in time vs.\ broad in
frequency'' do these two illustrate?
\end{exercise}

\begin{exercise}
By hand, find the $8$-point DFT of $x[n] = \cos(2\pi \cdot 3 n / 8)$. Which two
bins are nonzero and what are their values? Then do the same for
$x[n] = \sin(2\pi \cdot 3 n/8)$ and explain why the nonzero values are now
imaginary, connecting the result to the phase of the tone.
\end{exercise}

\begin{exercise}
You sample at $f_s = 1000$\,Hz and want to distinguish two tones $2$\,Hz apart.
(a) What is the minimum record length in seconds? (b) The minimum number of
samples $N$? (c) If instead you only had $N=256$ samples at this rate, what is your
bin spacing, and could you separate the two tones? Explain using
$\Delta f = 1/(NT_s)$.
\end{exercise}

\begin{exercise}
A tone of $4.5$ cycles per $N=16$ block is fed to a DFT. (a) Sketch the periodic
extension and mark the discontinuity at the seam. (b) Explain why no single bin
captures it. (c) Will multiplying by a Hann window before the DFT increase or
decrease the leaked side-lobes, and what does it cost you? Answer in terms of the
seam discontinuity.
\end{exercise}

\begin{exercise}
Confirm Parseval's identity for the $8$-point tone $x = (1,0,-1,0,1,0,-1,0)$ of
Example~\ref{ch08:ex-tone}. Compute $\sum_n x[n]^2$ directly, compute
$\frac{1}{N}\sum_k \abs{X[k]}^2$ from the two nonzero bins $X[2]=X[6]=4$, and show
they agree.
\end{exercise}

\begin{exercise}
Show that for any real input $x[n]$, the DFT satisfies $X[N-k] =
\overline{X[k]}$. (Hint: write out $X[N-k]$, use $e^{-\jj 2\pi (N-k)n/N} =
e^{\jj 2\pi k n/N}$, and recognize the conjugate.) Conclude that an $N$-point DFT
of real data has only $N/2 + 1$ independent complex values, and check that the
real-number count balances.
\end{exercise}

\begin{exercise}
Zero-padding. Take $x = (1,2,3,4)$ and append four zeros to make an $8$-point
input. Without grinding all eight sums, argue what zero-padding does to the
\emph{bin spacing} and to the \emph{frequency resolution} (are they the same
thing?). Then in NumPy compare \texttt{np.fft.fft(x)} with
\texttt{np.fft.fft(x, 8)} and describe how the second interpolates the first.
\end{exercise}

\begin{exercise}
Write out the $\mat{F}_8$ matrix's row $k=4$ explicitly and explain why it is the
``Nyquist row''. What time-domain signal does it perfectly match, and why is it the
fastest oscillation an $8$-sample block can represent?
\end{exercise}

\begin{exercise}
(IMU mini-project, conceptual.) You buffer $N=256$ samples of one accelerometer
axis at $f_s = 200$\,Hz and window with a Hann window. (a) Give the formula for the
center frequency of bin $k$. (b) Which bins cover the $1$--$5$\,Hz ``gesture''
band? (c) Explain in two or three sentences why summing $\abs{X[k]}^2$ over that
band is a more robust gesture feature than any single bin's magnitude. (d) State
one thing this wrist-IMU spectrum \emph{cannot} tell you about the hand, and which
sensor you would add to get it.
\end{exercise}

\begin{exercise}
(Matrix structure.) Verify numerically that $\mat{F}_4^{*}\mat{F}_4 = 4\mat{I}$ by
computing two off-diagonal inner products of rows and showing they vanish, plus one
diagonal inner product showing it equals $4$. Relate this to why
$\mat{F}_4^{-1} = \frac14 \mat{F}_4^{*}$.
\end{exercise}

\begin{exercise}
(Leakage experiment.) In NumPy, build $x[n] = \cos(2\pi f_0 n / f_s)$ for
$f_s = 256$, $N = 256$, and sweep $f_0$ from $10.0$ to $10.5$\,Hz in steps of
$0.1$. For each, plot $\abs{X[k]}$ near bin $10$. Describe how the single spike
turns into a two-bin split and a skirt as $f_0$ moves off-bin, and confirm the
peak height drops (scalloping loss).
\end{exercise}
